\documentclass[10pt,a4paper]{amsart}
\usepackage{amsmath,amssymb,amsthm,graphicx}
\usepackage{verbatim}
\newenvironment{reference}{\comment}{\endcomment}
%\newenvironment{reference}{}{}
\newenvironment{slogan}{\comment}{\endcomment}
\newenvironment{history}{\comment}{\endcomment}

% For commutative diagrams we use Xy-pic
\usepackage[all]{xy}

% We use 2cell for 2-commutative diagrams.
\xyoption{2cell}
% Theorem environments.
%
\theoremstyle{plain}
\newtheorem{theorem}[subsection]{Theorem}
\newtheorem{proposition}[subsection]{Proposition}
\newtheorem{lemma}[subsection]{Lemma}

\theoremstyle{definition}
\newtheorem{definition}[subsection]{Definition}
\newtheorem{example}[subsection]{Example}
\newtheorem{exercise}[subsection]{Exercise}
\newtheorem{situation}[subsection]{Situation}

\theoremstyle{remark}
\newtheorem{remark}[subsection]{Remark}
\newtheorem{remarks}[subsection]{Remarks}

\numberwithin{equation}{subsection}

% Macros
%
\def\lim{\mathop{\mathrm{lim}}\nolimits}
\def\colim{\mathop{\mathrm{colim}}\nolimits}
\def\Spec{\mathop{\mathrm{Spec}}}
\def\Hom{\mathop{\mathrm{Hom}}\nolimits}
\def\Ext{\mathop{\mathrm{Ext}}\nolimits}
\def\SheafHom{\mathop{\mathcal{H}\!\mathit{om}}\nolimits}
\def\SheafExt{\mathop{\mathcal{E}\!\mathit{xt}}\nolimits}
\def\Sch{\mathit{Sch}}
\def\Mor{\mathop{\mathrm{Mor}}\nolimits}
\def\Ob{\mathop{\mathrm{Ob}}\nolimits}
\def\Sh{\mathop{\mathit{Sh}}\nolimits}
\def\NL{\mathop{N\!L}\nolimits}
\def\CH{\mathop{\mathrm{CH}}\nolimits}
\def\proetale{{pro\text{-}\acute{e}tale}}
\def\etale{{\acute{e}tale}}
\def\QCoh{\mathit{QCoh}}
\def\Ker{\mathop{\mathrm{Ker}}}
\def\Im{\mathop{\mathrm{Im}}}
\def\Coker{\mathop{\mathrm{Coker}}}
\def\Coim{\mathop{\mathrm{Coim}}}

% Boxtimes
%
\DeclareMathSymbol{\boxtimes}{\mathbin}{AMSa}{"02}

%
% Macros for moduli stacks/spaces
%
\def\QCohstack{\mathcal{QC}\!\mathit{oh}}
\def\Cohstack{\mathcal{C}\!\mathit{oh}}
\def\Spacesstack{\mathcal{S}\!\mathit{paces}}
\def\Quotfunctor{\mathrm{Quot}}
\def\Hilbfunctor{\mathrm{Hilb}}
\def\Curvesstack{\mathcal{C}\!\mathit{urves}}
\def\Polarizedstack{\mathcal{P}\!\mathit{olarized}}
\def\Complexesstack{\mathcal{C}\!\mathit{omplexes}}
% \Pic is the operator that assigns to X its picard group, usage \Pic(X)
% \Picardstack_{X/B} denotes the Picard stack of X over B
% \Picardfunctor_{X/B} denotes the Picard functor of X over B
\def\Pic{\mathop{\mathrm{Pic}}\nolimits}
\def\Picardstack{\mathcal{P}\!\mathit{ic}}
\def\Picardfunctor{\mathrm{Pic}}
\def\Deformationcategory{\mathcal{D}\!\mathit{ef}}
\usepackage{cleveref}
\crefname{originalremark}{original remark}{original remarks}
\Crefname{originalremark}{Original remark}{Original remarks}
\makeatletter
\newlabel{remark-differential-object-selfmap}{{0AMK}{source}{}{sourcepointer-1}{}}
\newlabel{remark-differential-object-selfmap@cref}{{[originalremark][0][]0AMK}{[1][0][]source}{}{}{}}
\makeatother
\title{The spectral sequence associated to a filtered differential object}
\author{The Stacks Project Authors}
\date{stacks-project:homology.tex:local-snapshot-20261012; GNU FDL 1.2 or later}
\begin{document}
\maketitle
\noindent\textit{Editorial note.} All delivered mathematical text is the author's own, copied from the pinned Stacks Project chapter file homology.tex; no statement or proof step was rewritten, reordered or removed.\par\medskip
\noindent\textit{Editorial note.} The proof is the author's complete argument. It contains the author's own reading recommendation 'we strongly suggest the reader skip this proof', which by LANE\_BRIEF 5 is a suggestion to the reader and not a refusal to prove; the general construction is then carried out in full inside the counted span.\par\medskip
\noindent\textit{Editorial note.} Every labelled reference occurring inside the counted spans is carried in assembly.external\_references as a non-load-bearing navigation pointer (same-chapter ones at their definition site, cross-chapter ones at the invocation inside the counted span), so no load-bearing mathematics is replaced by a stub. No bibliography citation occurs inside the counted spans.\par\medskip
\noindent\textit{Editorial note.} The same-name check for this card scanned every staged Stacks chapter of this repository -- all sources/stacks-*/pinned-original-source.tex files, the card's own chapter excluded -- and the labels found are recorded in dedup.same\_name\_cross\_chapter on the card.\par\medskip
\noindent\textit{Editorial note.} The preamble is the chapter's own native preamble excerpt already used for this source: the comment-type reference, slogan and history environments (native preamble lines 8-13), the Xy-pic setup (native preamble lines 14-18) and the native theorem declarations with the mathematical macros (native preamble lines 154-217, the actual last line of the file).\par\medskip
\noindent\textit{Allowed standard prerequisites (editorial).}\begin{itemize}
\item Spectral sequences from an injective self-map of a differential object. Same chapter, Remark remark-differential-object-selfmap. Supplies the construction of a spectral sequence from an injective morphism of differential objects, applied with the shift [1], and the descriptions of the boundaries and cocycles used at the end.
\item Filtered differential objects and the graded category. Same chapter, the definitions of a filtered differential object, of the graded pieces gr\textasciicircum{}p and of the category Gr(A). Supplies the objects and the differential whose graded pieces are identified with the first two pages.
\item The countable direct sum case. Same chapter, the discussion immediately preceding the lemma. The opening sentence of the proof reduces to it in the case where the abelian category has exact countable direct sums.
\end{itemize}\medskip
\ifdefined\hypertarget\hypertarget{sourcepointer-1}{}\fi
\noindent\textit{Source reference.} Original remark 0AMK: sources/stacks-homology/pinned-original-source.tex, source lines 5811--5811. Navigation pointer invoked at proof line 5993; the counted span contains the pointer only and the cited result is not reproduced inside it, so the pointer is registered as a non-load-bearing navigation stub in the established form for Stacks material. No bibliography citation is involved.\par\medskip
\begin{lemma}
\label{lemma-spectral-sequence-filtered-differential}
Let $\mathcal{A}$ be an abelian category. Let $(K, F, d)$ be a
filtered differential object of $\mathcal{A}$. There is a
spectral sequence $(E_r, d_r)_{r \geq 0}$ in $\text{Gr}(\mathcal{A})$
associated to $(K, F, d)$ such that $d_r : E_r \to E_r[r]$
for all $r$ and such that the graded pieces
$E_r^p$ and maps $d_r^p : E_r^p \to E_r^{p + r}$
are as given above. Furthermore, $E_0^p = \text{gr}^p K$,
$d_0^p = \text{gr}^p(d)$, and $E_1^p = H(\text{gr}^pK, \text{gr}^p(d))$.
\end{lemma}

\begin{proof}
If $\mathcal{A}$ has countable direct sums and if countable direct
sums are exact, then this follows from the discussion above.
In general, we proceed as follows; we strongly suggest the reader
skip this proof. Consider the object $A = (F^{p + 1}K)$ of
$\text{Gr}(\mathcal{A})$, i.e., we put $F^{p + 1}K$ in degree $p$
(the funny shift in numbering to get numbering correct later on).
We endow it with a differential $d$ by using $d$ on each component.
Then $(A, d)$ is a differential object of $\text{Gr}(\mathcal{A})$.
Consider the map
$$
\alpha : A \to A[-1]
$$
which is given in degree $p$ by the inclusions $F^{p + 1}A \to F^pA$.
This is clearly an injective morphism of differential objects
$\alpha : (A, d) \to (A, d)[-1]$. Hence, we can apply
Remark \ref{remark-differential-object-selfmap}
with $S = \text{id}$ and $T = [1]$.
The corresponding spectral sequence $(E_r, d_r)_{r \geq 0}$
in $\text{Gr}(\mathcal{A})$ is the spectral sequence we are looking
for. Let us unwind the definitions a bit.
First of all we have $E_r = (E_r^p)$ is an object of $\text{Gr}(\mathcal{A})$.
Then, since $T^rS = [r]$ we have $d_r : E_r \to E_r[r]$ which means that
$d_r^p : E_r^p \to E_r^{p + r}$.

\medskip\noindent
To see that the description of the graded pieces hold, we argue
as above. Namely, first we have $E_0 = \Coker(\alpha : A \to A[-1])$
and by our choice of numbering above this gives
$E_0^p = \text{gr}^pK$. The first differential is given by
$d_0^p = \text{gr}^pd : E_0^p \to E_0^p$.
Next, the description of the boundaries $B_r$ and the cocycles $Z_r$
in Remark \ref{remark-differential-object-selfmap}
translates into a straightforward manner into the formulae
for $Z_r^p$ and $B_r^p$ given above.
\end{proof}

\end{document}
